\chapter{Justificación}
\label{chap:justificacion}

La prevalencia del comportamiento sedentario (CS) en la sociedad contemporánea es un fenómeno alarmante que afecta significativamente la salud y la calidad de vida de las personas \citep{Guthold2020,Swinburn2019Syndemic}. La Organización Mundial de la Salud (OMS) ha identificado el CS como uno de los principales factores de riesgo para enfermedades no transmisibles, tales como enfermedades cardiovasculares, diabetes tipo 2 y ciertos tipos de cáncer \citep{DiPietro2020Advancing,Murray2020GBD,Okunogbe2021Economic,WHO2009GlobalRisks}. Además, el sedentarismo tiene un impacto negativo en la salud mental y genera una disminución en la calidad de vida relacionada con la salud \citep{DiPietro2020Advancing,Guthold2020,Swinburn2019Syndemic}.

Los métodos tradicionales de evaluación del comportamiento sedentario presentan una dicotomía metodológica: los instrumentos subjetivos (cuestionarios de autoinforme) son escalables pero adolecen de sesgos de memoria y deseabilidad social, mientras que las técnicas objetivas de referencia (calorimetría indirecta, actigrafía de investigación) se restringen a entornos de laboratorio que no capturan la complejidad del comportamiento en condiciones de vida libre \citep{Healy2024,Prince2008,Migueles2022GRANADA}. Esta limitación metodológica fundamenta la necesidad de enfoques innovadores que integren objetividad (mediciones continuas mediante sensores), validez ecológica (recopilación en contextos reales) e interpretabilidad (clasificaciones auditables clínicamente).

La presente investigación abordó este vacío mediante el desarrollo y validación de un sistema de inferencia difusa tipo Mamdani que clasifica el comportamiento sedentario a partir de datos biométricos multivariados obtenidos de dispositivos portátiles de consumo (Apple Watch) bajo el paradigma \textit{Bring Your Own Device} (BYOD). Este enfoque metodológico integró tres componentes innovadores: (1) normalización fisiológica individualizada de variables (Actividad Relativa, Superávit Calórico Basal, Delta Cardíaco, HRV-SDNN), (2) establecimiento de una verdad operativa mediante clustering no supervisado (evitando umbrales heurísticos predefinidos), y (3) validación rigurosa mediante Leave-One-User-Out para garantizar generalización inter-sujeto.

La relevancia científica de este trabajo reside en su contribución al campo emergente de la inteligencia artificial explicable aplicada a salud digital. A diferencia de modelos de caja negra (redes neuronales profundas), el sistema difuso propuesto es inherentemente interpretable mediante reglas lingüísticas auditables, una característica crítica para la adopción clínica donde la transparencia algorítmica es un imperativo ético. Además, la demostración empírica de que el paradigma BYOD puede generar datos de calidad investigativa sin comprometer la integridad de las mediciones representa un avance significativo hacia la democratización de la investigación en salud, permitiendo estudios longitudinales a gran escala sin requerir equipamiento especializado costoso.

Desde la perspectiva de salud pública, esta investigación aporta evidencia sobre la factibilidad de sistemas de vigilancia del sedentarismo basados en datos pasivamente capturados por wearables de consumo masivo. Los resultados obtenidos (F1-Score LOOU = 0.780, validado contra clustering no supervisado) demuestran que es posible clasificar con alta fiabilidad los patrones de comportamiento sedentario en condiciones de vida libre, sentando las bases para el diseño de intervenciones digitales personalizadas y programas de monitoreo comunitario escalables. En un contexto nacional donde la prevalencia de sedentarismo alcanza el 40\% en adultos según ENSANUT 2022 \citep{Shamah-Levy2023ENSANUT}, la disponibilidad de herramientas objetivas, escalables e interpretables para su evaluación constituye una contribución estratégica para el cumplimiento de la meta 3.4 de los Objetivos de Desarrollo Sostenible (reducir en un tercio la mortalidad prematura por ENT para 2030).
